\section{From $A_\infty$ chiral identities to vertex Poisson axioms: full verification}
\label{app:Ainfty-to-PVA-axioms}

This appendix provides the complete, explicit verification that the cohomology $H = H^\bullet(A,Q)$ of an $A_\infty$ chiral algebra $(\A, \{m_k\}, Q)$ satisfying Hypotheses~(H1)--(H4) carries the structure of a $(-1)$-shifted Poisson vertex algebra. We verify each PVA axiom individually, with all signs tracked.

\medskip

\noindent\textbf{Standing notation.} Throughout, $[a],[b],[c]\in H$ denote cohomology classes with representatives $a,b,c\in \A$. The cohomological degree of $a$ is written $\degree{a}$. The product and bracket are defined from the regular/singular decomposition of $m_2$ as in Definition~\ref{def:product} and Definition~\ref{def:lambda-bracket}:
\begin{align}
\label{eq:app-product-recall}
[a]\cdot [b] &:= \Bigl[\,\tfrac{1}{2}\bigl(m_2^{\mathrm{reg}}(a,b) + (-1)^{\degree{a}\degree{b}}\,m_2^{\mathrm{reg}}(b,a)\bigr)\,\Bigr]\Big|_{\lambda=0},\\[4pt]
\label{eq:app-bracket-recall}
\{[a]{}_\lambda [b]\} &:= \Bigl[\,\tfrac{1}{2}\bigl(m_2^{\mathrm{sing}}(a,b;\lambda) - (-1)^{(\degree{a}+1)(\degree{b}+1)}\, e^{-(\lambda+\partial)\partial_\lambda}\, m_2^{\mathrm{sing}}(b,a;\lambda)\bigr)\,\Bigr].
\end{align}
Here $m_2^{\mathrm{reg}}(a,b) \in A[[\lambda]]$ and $m_2^{\mathrm{sing}}(a,b) \in \lambda^{-1}A[[\lambda^{-1}]]$ are the regular and singular parts (Definition~\ref{def:Laurent-projection}), $\partial$ is holomorphic translation, and $e^{-(\lambda+\partial)\partial_\lambda}$ implements the shift $\lambda \mapsto -\lambda - \partial$.

\subsection{(PVA1) Commutativity of the product}

\begin{proposition}[Commutativity]
\label{prop:app-commutativity}
For all $[a],[b]\in H$, the product satisfies $[a]\cdot [b] = (-1)^{\degree{a}\degree{b}}\,[b]\cdot [a]$.
\end{proposition}

\begin{proof}
By definition \eqref{eq:app-product-recall},
\[
[a]\cdot[b] = \Bigl[\,\tfrac{1}{2}\bigl(m_2^{\mathrm{reg}}(a,b) + (-1)^{\degree{a}\degree{b}}\,m_2^{\mathrm{reg}}(b,a)\bigr)\,\Bigr]\Big|_{\lambda=0}.
\]
Computing $(-1)^{\degree{a}\degree{b}}\,[b]\cdot[a]$:
\begin{align*}
(-1)^{\degree{a}\degree{b}}\,[b]\cdot[a]
&= (-1)^{\degree{a}\degree{b}}\Bigl[\,\tfrac{1}{2}\bigl(m_2^{\mathrm{reg}}(b,a) + (-1)^{\degree{b}\degree{a}}\,m_2^{\mathrm{reg}}(a,b)\bigr)\,\Bigr]\Big|_{\lambda=0}\\
&= \Bigl[\,\tfrac{1}{2}\bigl((-1)^{\degree{a}\degree{b}}\,m_2^{\mathrm{reg}}(b,a) + m_2^{\mathrm{reg}}(a,b)\bigr)\,\Bigr]\Big|_{\lambda=0}\\
&= [a]\cdot[b].
\end{align*}
The crucial point is that the graded symmetrization $\Sym$ in the definition of the product builds commutativity in \emph{by construction}: writing
\[
\Sym\bigl(m_2^{\mathrm{reg}}(a,b),\,m_2^{\mathrm{reg}}(b,a)\bigr) = \tfrac{1}{2}\bigl(m_2^{\mathrm{reg}}(a,b) + (-1)^{\degree{a}\degree{b}}\,m_2^{\mathrm{reg}}(b,a)\bigr),
\]
exchanging $a\leftrightarrow b$ multiplies the expression by $(-1)^{\degree{a}\degree{b}}$, which is exactly the graded commutativity condition.

\smallskip

\noindent\textbf{Remark on well-definedness.} One must also verify that the symmetrization is compatible with passing to cohomology, i.e., that the antisymmetric part $m_2^{\mathrm{reg}}(a,b) - (-1)^{\degree{a}\degree{b}}\,m_2^{\mathrm{reg}}(b,a)$ is $Q$-exact. This follows from the $\FM_2(\C)$ involution argument: the map $\sigma: (z_1,z_2)\mapsto (z_2,z_1)$ acts on $\FM_2(\C)$, and the weight form transforms as
\[
\sigma^*\omega_2(a,b) = (-1)^{\degree{a}\degree{b}}\,\omega_2(b,a).
\]
The involution $\sigma$ is connected to the identity through paths in $\Conf_2(\R)$ (using the topological direction), and this path provides a chain homotopy $h$ with
\[
m_2^{\mathrm{reg}}(a,b) - (-1)^{\degree{a}\degree{b}}\,m_2^{\mathrm{reg}}(b,a) = Q\,h(a,b) + h(Qa,b) + (-1)^{\degree{a}} h(a,Qb).
\]
On $Q$-closed representatives with $Qa = Qb = 0$, the right side is $Q$-exact, so the symmetrization is exact in cohomology.
\end{proof}

\subsection{(PVA2) Sesquilinearity of the $\lambda$-bracket}

\begin{proposition}[Sesquilinearity]
\label{prop:app-sesqui}
The $\lambda$-bracket satisfies:
\begin{align}
\label{eq:app-sesqui-L}
\{\partial [a]{}_\lambda [b]\} &= -\lambda\,\{[a]{}_\lambda [b]\},\\[4pt]
\label{eq:app-sesqui-R}
\{[a]{}_\lambda \partial [b]\} &= (\lambda + \partial)\,\{[a]{}_\lambda [b]\}.
\end{align}
\end{proposition}

\begin{proof}
\textbf{Step 1: Left sesquilinearity \eqref{eq:app-sesqui-L}.}
The chain-level sesquilinearity (equation~\eqref{eq:sesqui-left} of Definition~\ref{def:sesquilinearity}) states
\[
m_2(\partial a,\, b) = -\lambda \cdot m_2(a,b).
\]
This is an identity of Laurent series in $\lambda$. We project to the singular part. Since multiplication by $-\lambda$ maps the singular part to itself (concretely, if $m_2^{\mathrm{sing}}(a,b) = \sum_{n\ge 1} c_n/\lambda^n$, then
\[
-\lambda \cdot m_2^{\mathrm{sing}}(a,b) = -\lambda \cdot \sum_{n\ge 1}\frac{c_n}{\lambda^n} = -c_1 - \sum_{n\ge 2}\frac{c_n}{\lambda^{n-1}},
\]
while $-\lambda \cdot m_2^{\mathrm{reg}}(a,b) = -\lambda \cdot \sum_{n\ge 0} d_n \lambda^n$ has no singular part (all terms have non-negative powers of $\lambda$)), we conclude
\[
m_2^{\mathrm{sing}}(\partial a,\, b) = \bigl[-\lambda \cdot m_2(a,b)\bigr]_{\mathrm{sing}} = -\lambda \cdot m_2^{\mathrm{sing}}(a,b).
\]

\noindent\textbf{Detail on the cross-term.} Strictly, $-\lambda \cdot m_2^{\mathrm{sing}}(a,b)$ has a regular contribution: $-\lambda \cdot (c_1/\lambda) = -c_1 \in A[[\lambda]]$. Meanwhile $-\lambda \cdot m_2^{\mathrm{reg}}(a,b)$ contributes $-\lambda d_0 - \lambda^2 d_1 - \cdots$, which is purely regular. The full singular part of $-\lambda \cdot m_2(a,b)$ is therefore $-\sum_{n\ge 2} c_n/\lambda^{n-1} = -\lambda \cdot \sum_{n\ge 2} c_n/\lambda^n$.

However, this subtlety does not affect the $\lambda$-bracket, because the bracket is defined via the \emph{antisymmetrized} singular part~\eqref{eq:app-bracket-recall}. Both the $m_2^{\mathrm{sing}}(\partial a, b)$ term and the $e^{-(\lambda+\partial)\partial_\lambda}\,m_2^{\mathrm{sing}}(b,\partial a)$ term in the antisymmetrization transform consistently under the chain-level identity. Explicitly, using sesquilinearity in the second slot (equation~\eqref{eq:sesqui-right}):
\[
m_2(b,\partial a) = (\lambda + \partial)\,m_2(b,a),
\]
so that $m_2^{\mathrm{sing}}(b,\partial a) = [(\lambda+\partial)\,m_2(b,a)]_{\mathrm{sing}}$. Applying $e^{-(\lambda+\partial)\partial_\lambda}$ and collecting terms, the full antisymmetrized bracket transforms by $-\lambda$, giving
\[
\{\partial [a]{}_\lambda [b]\} = -\lambda\,\{[a]{}_\lambda [b]\}.
\]

\textbf{Step 2: Right sesquilinearity \eqref{eq:app-sesqui-R}.}
The chain-level identity \eqref{eq:sesqui-right} gives
\[
m_2(a,\, \partial b) = (\lambda + \partial)\,m_2(a,b).
\]
Projecting to the singular part: the operator $(\lambda + \partial)$ decomposes as $\lambda \cdot (\text{id}) + \partial$. Since $\partial$ acts on the coefficient algebra $A$ and does not involve $\lambda$, it commutes with the singular/regular decomposition:
\[
[\partial \cdot F(\lambda)]_{\mathrm{sing}} = \partial \cdot [F(\lambda)]_{\mathrm{sing}}.
\]
For the $\lambda$-multiplication piece, we again use that $\lambda$ shifts the power of $\lambda$ by $+1$, potentially converting one singular term ($c_1/\lambda$) to a regular constant. Tracking through both slots of the antisymmetrization \eqref{eq:app-bracket-recall}:

For the first slot, $m_2^{\mathrm{sing}}(a, \partial b) = [(\lambda + \partial)\,m_2(a,b)]_{\mathrm{sing}}$.

For the second slot, we need $m_2^{\mathrm{sing}}(\partial b, a)$. By \eqref{eq:sesqui-left}, $m_2(\partial b, a) = -\lambda \cdot m_2(b,a)$. So $m_2^{\mathrm{sing}}(\partial b, a) = [-\lambda \cdot m_2(b,a)]_{\mathrm{sing}}$.

Applying the exponential shift $e^{-(\lambda+\partial)\partial_\lambda}$ to the second slot and combining with the first slot, the net effect on the antisymmetrized bracket is multiplication by $(\lambda + \partial)$:
\[
\{[a]{}_\lambda \partial [b]\} = (\lambda + \partial)\,\{[a]{}_\lambda [b]\}.
\]

\textbf{Step 3: Descent to cohomology.}
Since $[\partial, Q] = 0$ by hypothesis and $\lambda$ is a formal variable commuting with $Q$, both sesquilinearity relations hold at the chain level and descend to cohomology without modification.
\end{proof}

\subsection{(PVA3) Skew-symmetry}

\begin{proposition}[Skew-symmetry]
\label{prop:app-skew}
The $\lambda$-bracket satisfies
\begin{equation}
\label{eq:app-skew}
\{[a]{}_\lambda [b]\}
= -(-1)^{\degree{a}\degree{b}}\,
\{[b]{}_{-\lambda-\partial}\,[a]\}.
\end{equation}
\end{proposition}

\begin{proof}
Let \(\zeta=z_1-z_2\).  The exchange cylinder traverses the
holomorphic half-boundary arc \(\zeta(s)=2r e^{\pi i s}\) and
transposes the topological \(E_1\)-ordering.  Pulling the two-point
logarithmic kernel to this oriented cylinder and applying Cartan--Stokes
gives, after singular projection,
\[
m_2^{\mathrm{sing}}(a,b;\zeta)
+(-1)^{\degree{a}\degree{b}}\tau^\ast
m_2^{\mathrm{sing}}(b,a;\zeta)
=QH_{\mathrm{ex}}^{\mathrm{sing}}(a,b;\zeta)
\]
up to the standard \(Q\)-exact input-correction terms.  The operator
\(\tau^\ast\) exchanges the two inputs and sends \(\zeta\mapsto-\zeta\).
Since \(d\log(-\zeta)=d\log\zeta\), the logarithmic form supplies no
additional sign; the minus sign in cohomology is the oriented endpoint
sign of the Stokes identity.

On \(Q\)-cohomology the right-hand side vanishes.  Applying the
polynomial Borel transform and translation covariance converts
\(\tau^\ast\) into the single vertex-algebra substitution
\(\lambda\mapsto-\lambda-\partial\).  Hence
\[
\{[a]{}_\lambda [b]\}
=-(-1)^{\degree{a}\degree{b}}\,
\{[b]{}_{-\lambda-\partial}[a]\}.
\]
\end{proof}

\subsection{(PVA4) Vacuum axiom}

\begin{proposition}[Vacuum]
\label{prop:app-vacuum}
Let $\mathbf{1}\in A^0$ be the strict unit with $Q\mathbf{1} = 0$ and $\partial\mathbf{1} = 0$. Then $[\mathbf{1}]\in H^0$ satisfies:
\begin{align}
\label{eq:app-vac-bracket}
\{[\mathbf{1}]{}_\lambda [a]\} &= 0,
\qquad
\{[a]{}_\lambda[\mathbf{1}]\}=0,\\[3pt]
\label{eq:app-vac-product}
[\mathbf{1}]\cdot [a] &= [a]\cdot[\mathbf{1}]=[a],
\qquad
\partial[\mathbf 1]=0.
\end{align}
\end{proposition}

\begin{proof}
\textbf{Step 1: Bracket vanishing \eqref{eq:app-vac-bracket}.}
The strict unit axiom \eqref{eq:unit-m2}, after the bar
desuspension used in the product convention, gives ordinary unit
identities for the unsuspended binary product:
\[
\mu_2^{\mathrm{reg}}(\mathbf 1,a;0)=a,\qquad
\mu_2^{\mathrm{reg}}(a,\mathbf 1;0)=a.
\]
The corresponding binary kernels are constant in the collision
parameter and have no singular part:
\[
m_2^{\mathrm{sing}}(\mathbf{1},a) = 0,
\qquad
m_2^{\mathrm{sing}}(a,\mathbf{1}) = 0.
\]
Since the $\lambda$-bracket is the divided-power Borel transform of
the singular part, \(\{[\mathbf 1]{}_\lambda[a]\}=0\) and
\(\{[a]{}_\lambda[\mathbf 1]\}=0\).

\textbf{Step 2: Unit for the product \eqref{eq:app-vac-product}.}
The descended product is represented by the unsuspended regular
constant term
\[
\mu(a,b):=\mu_2^{\mathrm{reg}}(a,b;0),
\]
and graded commutativity is proved separately by the exchange
homotopy.  Therefore
\[
[\mathbf 1]\cdot[a]=[\mu(\mathbf 1,a)]=[a],
\qquad
[a]\cdot[\mathbf 1]=[\mu(a,\mathbf 1)]=[a].
\]

\textbf{Step 3: Translation of the vacuum.}
By hypothesis, $\partial\mathbf{1} = 0$ and $Q\mathbf{1} = 0$, so $[\partial\mathbf{1}] = 0$ in $H$. This is consistent with sesquilinearity: $\{\partial\mathbf{1}{}_\lambda a\} = -\lambda\cdot\{\mathbf{1}{}_\lambda a\} = 0$.

\textbf{Step 4: Higher operations annihilate the vacuum.}
The strict unit axiom \eqref{eq:unit-higher} gives $m_k(\ldots,\mathbf{1},\ldots) = 0$ for all $k\neq 2$ whenever a unit is inserted. Combined with the vanishing of both singular binary projections with a unit insertion, no hidden higher boundary term modifies the vacuum identities.
\end{proof}

\subsection{(PVA5) Jacobi identity}

\begin{proposition}[Jacobi identity]
\label{prop:app-Jacobi}
For all $[a],[b],[c]\in H$, the $\lambda$-bracket satisfies
\begin{equation}
\label{eq:app-Jacobi}
\{[a]{}_\lambda \{[b]{}_\mu [c]\}\} - (-1)^{(\degree{a}+1)(\degree{b}+1)}\,\{[b]{}_\mu \{[a]{}_\lambda [c]\}\} = \{\{[a]{}_\lambda [b]\}{}_{\lambda+\mu}\, [c]\}.
\end{equation}
\end{proposition}

The Jacobi identity is the deepest of the PVA axioms and is proved in full detail in the main text (Theorem~\ref{thm:Jacobi}). Here we summarize the argument and explain the sign structure.

\begin{proof}[Proof (summary; see Theorem~\ref{thm:Jacobi} for full details)]
\textbf{Step 1: Source, the $n=3$ $A_\infty$ identity.}
The Stasheff identity at arity $n=3$ (equation~\eqref{eq:ainfty-relation-raw}) reads:
\begin{align*}
0 &= m_1(m_3(a,b,c)) + m_3(m_1(a),b,c) + (-1)^{\degree{a}}m_3(a,m_1(b),c) + (-1)^{\degree{a}+\degree{b}}m_3(a,b,m_1(c))\\
&\quad + m_2\bigl(m_2(a,b)\big|_{\Lambda_{12}},\, c;\, \Lambda_{12},\lambda_2\bigr) + (-1)^{\degree{a}}\,m_2\bigl(a,\, m_2(b,c)\big|_{\lambda_2};\, \lambda_1, \lambda_2\bigr),
\end{align*}
where $\Lambda_{12} = \lambda_1 + \lambda_2$ is the effective spectral parameter for the fused block $\{1,2\}$.

\textbf{Step 2: Pass to cohomology.}
On $Q$-closed representatives ($m_1(a) = Qa = 0$, etc.), the
total-collision term \(m_1(m_3(a,b,c))=Qm_3(a,b,c)\) is a local
\(Q\)-boundary in the arity-\(3\) Stokes identity.  This local
boundary is what is used in the Jacobi proof; it is not the separate
global statement that all higher operations vanish in cohomology.
After discarding this boundary term, the surviving relation in \(H\)
is:
\[
[m_2(m_2(a,b),c)] = [m_2(a,m_2(b,c))] \quad\text{in }H,
\]
which is the associativity of $m_2$ in cohomology.

\textbf{Step 3: Extract the Jacobi identity.}
Project the Stokes identity to the singular$\times$singular part (both
inner and outer binary channels singular).  Write the Jacobiator in the
normal form
\[
J_{\lambda,\mu}
=
\{a{}_\lambda \{b{}_\mu c\}\}
-
(-1)^{(\degree{a}+1)(\degree{b}+1)}
\{b{}_\mu \{a{}_\lambda c\}\}
-
\{\{a{}_\lambda b\}{}_{\lambda+\mu}c\}.
\]
The oriented pairwise boundary contributions of $\FM_3(\C)$ are:
\begin{itemize}
\item $D_{\{2,3\}}$: the composition $m_2(a,m_2(b,c))$ gives
 $\{a{}_\lambda \{b{}_\mu c\}\}$.
\item $D_{\{1,3\}}$: the non-consecutive dressed singular channel gives
 $-(-1)^{(\degree{a}+1)(\degree{b}+1)}\,\{b{}_\mu \{a{}_\lambda c\}\}$,
 the minus being the incidence orientation.
\item $D_{\{1,2\}}$: the composition $m_2(m_2(a,b),c)$ at spectral
 parameter $\Lambda_{12} = \lambda + \mu$ gives
 $-\{\{a{}_\lambda b\}{}_{\lambda+\mu}\,c\}$ in the displayed
 Jacobiator normalization.
\end{itemize}
The Arnold relation among the three logarithmic 1-forms
\[
d\log(z_1-z_2)\wedge d\log(z_2-z_3) + d\log(z_2-z_3)\wedge d\log(z_3-z_1) + d\log(z_3-z_1)\wedge d\log(z_1-z_2) = 0
\]
ensures that this signed boundary sum is a \(Q\)-boundary.  Hence
\(J_{\lambda,\mu}\) vanishes on cohomology, which is exactly
\eqref{eq:app-Jacobi}.

\textbf{Step 4: Sign analysis.}
The Koszul sign $(-1)^{(\degree{a}+1)(\degree{b}+1)}$ in \eqref{eq:app-Jacobi} differs from the naive sign $(-1)^{\degree{a}\degree{b}}$ by the shift $+1$ in each degree. This shift arises because the $\lambda$-bracket has degree $-1$ (it is a $(-1)$-shifted operation), so elements effectively carry shifted degrees $\degree{a}+1$ in the bracket. Explicitly:
\[
(-1)^{(\degree{a}+1)(\degree{b}+1)} = (-1)^{\degree{a}\degree{b}+\degree{a}+\degree{b}+1}.
\]
The $+1$ in the exponent provides the crucial ``super-Lie'' sign: in the unshifted setting, $(-1)^{\degree{a}\degree{b}}$ would give a graded Lie bracket, but the degree shift introduces an additional $(-1)^{\degree{a}+\degree{b}+1}$ factor.
\end{proof}

\subsection{(PVA6) Leibniz rule}

\begin{proposition}[Leibniz rule]
\label{prop:app-Leibniz}
For all $[a],[b],[c]\in H$:
\begin{equation}
\label{eq:app-Leibniz}
\{[a]{}_\lambda ([b]\cdot[c])\}
=
\{[a]{}_\lambda [b]\}\cdot [c]
+
(-1)^{(\degree{a}+1)\degree{b}}\,
[b]\cdot\{[a]{}_\lambda [c]\}.
\end{equation}
\end{proposition}

\begin{proof}
\textbf{Step 1: Source, the mixed regular--singular projection.}
Project the arity-$3$ Stokes identity on $\FM_3(\C)$ to the sector in
which exactly one binary channel is singular and the other binary
channel is the regular constant term.  Define the signed Leibniz
defect
\[
\mathcal L_\lambda(a;b,c)
:=
\{a{}_\lambda (b\cdot c)\}
-
\{a{}_\lambda b\}\cdot c
-
(-1)^{(\degree{a}+1)\degree{b}}\,
b\cdot\{a{}_\lambda c\}.
\]

\textbf{Step 2: The three mixed faces.}
The oriented pairwise boundary contributions are
\[
\mathfrak L_{23}\leadsto \{a{}_\lambda(b\cdot c)\},\qquad
\mathfrak L_{12}\leadsto -\,\{a{}_\lambda b\}\cdot c,\qquad
\mathfrak L_{13}\leadsto
-\,(-1)^{(\degree{a}+1)\degree{b}}\,
b\cdot\{a{}_\lambda c\}.
\]
The first term is obtained by taking the \(b,c\) collision in the
regular channel and the outer \(a\)-channel in the singular channel:
\[
\{[a]{}_\lambda ([b]\cdot [c])\}
=
\bigl[
m_2^{\mathrm{sing}}\bigl(a,\,
m_2^{\mathrm{reg}}(b,c)\big|_{\mu=0};\lambda\bigr)
\bigr].
\]
The \(a,b\) face gives the first right-hand term, with the minus sign
because we are computing the defect:
\[
\{[a]{}_\lambda [b]\}\cdot [c]
=
\bigl[
m_2^{\mathrm{reg}}\bigl(
m_2^{\mathrm{sing}}(a,b;\lambda),c;\mu\bigr)
\bigr]\Big|_{\mu=0}.
\]
The \(a,c\) face gives the second right-hand term; moving the
degree-\((-1)\) operator \(\{a{}_\lambda-\}\) across \(b\) produces the
shifted Poisson sign:
\[
(-1)^{(\degree{a}+1)\degree{b}}\,
[b]\cdot\{[a]{}_\lambda [c]\}
=
(-1)^{(\degree{a}+1)\degree{b}}
\bigl[
m_2^{\mathrm{reg}}\bigl(
b,m_2^{\mathrm{sing}}(a,c;\lambda);\mu\bigr)
\bigr]\Big|_{\mu=0}.
\]

\textbf{Step 3: Stokes and the absence of a Wick correction.}
The total boundary contribution vanishes by Stokes' theorem on $\FM_3(\C)$:
\[
0 = \sum_S \epsilon_{D_S}\int_{D_S}\Res_{D_S}(\omega_3),
\]
and the total-collision face is \(Q\)-exact.  Therefore
\(\mathcal L_\lambda(a;b,c)\in\operatorname{im}Q\), so the defect
vanishes after passing to \(H\).  This is exactly
\eqref{eq:app-Leibniz}.  No Wick-type integral term appears in the
classical Leibniz rule: the double singular bracket belongs to the
singular--singular Jacobi sector, not to the mixed regular--singular
sector.

\textbf{Step 4: Sign verification for the Koszul factor.}
The sign $(-1)^{(\degree{a}+1)\degree{b}}$ in the second term arises as follows. In the composition $m_2(b, \{a{}_\lambda c\})$, we must commute $b$ (of degree $\degree{b}$) past the bracket $\{a{}_\lambda -\}$ (of degree $\degree{a}+1-2 = \degree{a}-1$ as an operation, but the bracket output carries degree $\degree{a}+\degree{c}-1$, and the relevant sign is from commuting the bracket \emph{operation} past $b$). The shifted degree of $\{a{}_\lambda -\}$ is $\degree{a}+1$ (the $+1$ from the $(-1)$-shifted structure), giving the sign $(-1)^{(\degree{a}+1)\degree{b}}$.
\end{proof}

\subsection{Conclusion}

\begin{corollary}[Shifted PVA structure, restated]
Under Hypotheses~\textrm{(H1)--(H4)}, the cohomology $H = H^\bullet(\A,Q)$ with the product $\cdot$ (equation~\eqref{eq:app-product-recall}), $\lambda$-bracket $\{-{}_\lambda -\}$ (equation~\eqref{eq:app-bracket-recall}), translation $\partial$, and vacuum $[\mathbf{1}]$ is a $(-1)$-shifted Poisson vertex algebra: it satisfies \textrm{(PVA1)--(PVA6)} as verified in Propositions~\ref{prop:app-commutativity}--\ref{prop:app-Leibniz} above. This completes the proof of Corollary~\ref{cor:PVA-structure}.
\end{corollary}
